%	JASA LaTeX Sample File, Reprint Sample
%
%  Beginner Latex users should refer to their favorite online documentation
%  here is one from the TeX Users Group 
%	https://www.tug.org/twg/mactex/tutorials/ltxprimer-1.0.pdf
%
%  Useful FAQ from  https://journals.aps.org/revtex/revtex-faq
% 

%%%%%%% For Preprint
%% For manuscript, 12pt, one column style. Preprint is required for submission.

% \documentclass[preprint]{JASA}

%%%%% Preprint Options %%%%%
%% The track changes option allows you to mark changes
%% and will produce a list of changes, their line number
%% and page number at the end of the article.
%\documentclass[preprint,trackchanges]{JASA}

%% NumberedRefs is used for numbered bibliography and citations.
%% Default is Author-Year style.
% \documentclass[preprint,NumberedRefs]{JASA}

%%%%%%% For Reprint
%% For appearance of finished article; 2 columns, 10 pt fonts
%% Also used for estimating print page length.

 \documentclass[reprint]{JASA}

%%%%% Reprint Options %%%%%

%% NumberedRefs is used for numbered bibliography and citations.
%% Default is Author-Year style.
% \documentclass[reprint,NumberedRefs]{JASA}

%% TurnOnLineNumbers
%% Make lines be numbered in reprint style:
% \documentclass[reprint,TurnOnLineNumbers]{JASA}

\begin{document}

\title[JASA/Sample JASA Article]{Article title should be less than 17 words, no acronyms}
\author{Author One}
\author{Author Two}
\author{Author Three}
\affiliation{Department1,  University1, City, State ZipCode, Country}

\author{Author Four}
\email{author.four@university.edu}
\affiliation{Department2,  University2, City, State ZipCode, Country}
 
\author{Author Five}			
\altaffiliation{Also at: Another University, City, State ZipCode, Country.}
\affiliation{Department3,  University3, City, State ZipCode, Country}

%\preprint{Author, JASA}		
%  Used if you want want this message to appear in upper right corner of title page, for preprint only.

\date{\today} 

\begin{abstract}
Put your abstract here. Abstracts are limited to 200 words for
regular articles and 100 words for Letters to the Editor. Please no
personal pronouns, also please do not use the words ``new'' and/or
``novel'' in the abstract. An article usually includes an abstract, a
concise summary of the work covered at length in the main body of the
article.     
\end{abstract}

%% pacs numbers not used

\maketitle


\section{\label{sec:1} Introduction}
This sample document demonstrates the use of JASA in manuscripts 
prepared for submission to the Journal of the Acoustical Society of America. 

See JASA-TeXGuide.pdf, which is part of this package, for extensive
documentation on using commands for JASA.

You can compare the .tex version of this file with the resulting .pdf
version to give you an idea of what  commands are available and how
they work. At the top of the .tex file you'll find a listing of the
documentclass options, and an explanation of their results.
Some additional suggestions are included in the body of this
manuscript.  

  Beginner Latex users should refer to their favorite online documentation. A 
  useful place to start is the primer from the TeX Users Group 
  \url{https://www.tug.org/twg/mactex/tutorials/ltxprimer-1.0.pdf}


The paper is organized as follows: Section~\ref{sec:2} presents
initial information, while
Section~\ref{sec:3} presents examples of mathematical expressions.

{\it
[For these sample pages we have used only figsamp.jpg for
convenience]}

Here is some normal text that
show how the columns are formatted.
Here is some normal text that
show how the columns are formatted.
Here is some normal text that
show how the columns are formatted.
Here is some normal text that
show how the columns are formatted.
Here is some normal text that
show how the columns are formatted.
Here is some normal text that
show how the columns are formatted.


\section{\label{sec:2} Section Two}
An example of another first-level Section with following example text that refers to subsections using 
\verb+\ref{subsec:XXX}+ ...  EXAMPLE: Some background in
section~\ref{sec:2} and details  in subsection~\ref{subsec:2:1}. 

\subsection{\label{subsec:2:1} Sample subsection}

\subsubsection{Sample subsubsection\label{subsubsec:1}}





% Sample showing how to include figure, this is a floating figure

 
\begin{figure}
%% \reprintcolumnwidth is the same in preprint and reprint for
%% ease of use for authors:
\includegraphics[width=\reprintcolumnwidth]{figsamp}
\caption{\label{fig:FIG1}{Caption here.}}
\raggedright
Note: The only figure formats allowed are the following: 
.pdf, .ps, .eps, or .jpg.  (.tiff is not recommended for use with LaTeX)

Figure files must be named in this fashion:
Figure\#.xxx, where ``\#'' is the figure number and ``xxx'' is the file format
(Figure1.eps, Figure2.jpg, Figure3a.ps, Figure3b.ps, etc). 

\end{figure}

\paragraph{Sample paragraph}Here is text following the paragraph
heading.
Here is a figure reference: is shown in Fig.~\ref{fig:FIG1}.


\section{Inline and display math samples\label{sec:3}}
\subsection{Tip: How to keep line numbers from disappearing above
some math commands} If you use line numbering \\
(\verb+\documentclass[reprint,TurnOnLineNumbers]{JASA}+)\\
occasionally line numbers will disappear in the paragraph above
certain math commands, like \verb+\begin{align}+.
You can fix this by typing \verb+\linenomath+ before the math command,
which will allow the numbering to continue in the paragraph before the
math:
\linenomath
\begin{align} 
 2x - 5y &=  8 \\ 
 3x + 9y &=  -12
\end{align}
Another line here.

\subsection{\label{subsec:3:3} Math and equations $\alpha\beta\Delta\Gamma$}
Inline math may be typeset using the \verb+$+ delimiters.  (Authors:
Remember to surround your math with the \verb+$+ delimiters.  A missing dollar sign is a common cause for errors.)

Bold math
symbols may be achieved using the \verb+bm+ package and the
\verb+\bm{#1}+ command it supplies. For instance, a bold $\alpha$ can
be typeset as \verb+$\bm{\alpha}$+ giving $\bm{\alpha}$. Fraktur and
Blackboard (or open face or double struck) characters should be
typeset using the \verb+\mathfrak{#1}+ and \verb+\mathbb{#1}+ commands
respectively. Both are supplied by the \texttt{amssymb} package which
is called in JASA, so no need enter \verb+\includepackage{amssymb}+.
To see amssymb in action:
\verb+$\mathbb{R}$+ gives $\mathbb{R}$ and
\verb+$\mathfrak{G}$+ gives $\mathfrak{G}$.

In \LaTeX\ there are many different ways to display equations; a
few preferred ways are noted below. Displayed math will center by
default. 

Below we have numbered single-line equations; this is the most common
type of equation.
\begin{eqnarray}
\chi_+(p){\bf [}2|{\bf p}|(|{\bf p}|+p_z){\bf ]}^{-1/2}
\left(
\begin{array}{c}
|{\bf p}|+p_z\\
px+ip_y
\end{array}\right)\;,
\\
\left\{%
 1234567890abc123\alpha\beta\gamma\delta1234556\alpha\beta
 \frac{1\sum^{a}_{b}}{A^2}%
\right\}%
\label{eq:one}.
\end{eqnarray}
Note the open one in Eq.~(\ref{eq:one}).

Not all numbered equations will fit within a narrow column this
way. The equation number will move down automatically if it cannot fit
on the same line with a one-line equation.
\begin{equation}
\chi_+(p){\bf [}2|{\bf p}|(|{\bf p}|+p_z){\bf ]}^{-1/2}
\alpha\beta\gamma\delta123455678\alpha\beta\Gamma\Delta
 \frac{1\sum^{a}_{b}}{A^2}1234%
\end{equation}
% No blank lines after \end{equation}
When the \verb+\label{#1}+ command is used [ie. input for
Eq.~(\ref{eq:one})], the equation can be referred to in text without
knowing the equation number that \TeX\ will assign to it. Just
use \verb+\ref{#1}+, where \verb+#1+ is the same name that used in
the \verb+\label{#1}+ command.

Unnumbered single-line equations can be typeset
using the \verb+\[+, \verb+\]+ format:
\[g^+g^+ \rightarrow g^+g^+g^+g^+ \dots ~,~~q^+q^+\rightarrow
q^+g^+g^+ \dots ~. \]
Equations can be lettered with the
subequations environment:
\begin{subequations}
\begin{eqnarray}
A&=&mc, \label{eqa}
\\
B&=&mc^2, \label{eqb}
\\
C&\agt& mc^3. \label{eqc}
\end{eqnarray}
\end{subequations}
Referenced: Eqs.~(\ref{eqa}), (\ref{eqb}), and (\ref{eqc}).


\section{Floats, Figures and Tables}

 Figures and tables are typically ``floats'' which means that their
final position is determined by \LaTeX\ while the document is being
typeset. \LaTeX\ isn't always successful in placing floats
optimally.  Use the figure* environment to get a wide figure that spans the 
page in a two-column layout.

\subsection{\label{subsec:3:2} Tables using \tt\string\ruledtabular}
Tables generally should be surrounded with
\verb+\begin{ruledtabular}...\end{ruledtabular}+\\
This will guarantee that they are the width of the
page or column, and have two ruled lines at the top
and bottom of the table.

\verb+[ht]+ in the code below instructs \LaTeX\ to place the table
where it appears in type, if it will fit on the page;
otherwise put it on the top of the next page.

\subsubsection{Table notes}
Footnotes in a table are labeled a, b, c,
 etc.  They can
be  specified  by  using  the  \LaTeX\
\verb+\footnotemark[]+
and
\verb+\footnotetext[]+ commands.
The  footnotes  for  a  table  are  typeset  at  the
bottom  of  the  table,  rather  than  at  the  bottom  of  the
page or at the end of the references.  The arguments for
\verb+\footnotemark[]+
and
\verb+\footnotetext[]+
should be numbers
1, 2, \ldots\  The journal style will convert these to letters.

This system allows
multiple entries to refer to the same
footnote.   


\begin{table}[ht]
\caption{\label{tab:table1}A table using {\tt\string\ruledtabuler}.
Note that several entries share the same
footnote. Inspect the \LaTeX\ input for this table to see
exactly how it is done.}

\begin{ruledtabular}
\begin{tabular}{cccccccc}
 &$r_c$ (\AA)\footnotemark[1]&$r_0$ (\AA)&$\kappa r_0$&
 &$r_c$ (\AA) &$r_0$ (\AA)&$\kappa r_0$\\
\hline
Cu& 0.800 & 14.10 & 2.550 &Sn\footnotemark[1]
& 0.680 & 1.870 & 3.700 \\
Ag& 0.990 & 15.90 & 2.710 &Pb\footnotemark[2]
& 0.450 & 1.930 & 3.760 \\
Au& 1.150 & 15.90 & 2.710 &Ca\footnotemark[3]
& 0.750 & 2.170 & 3.560 
\end{tabular}
\end{ruledtabular}
\footnotetext[1]{Here's the first.}
\footnotetext[2]{Here's the second.}
\footnotetext[3]{Here's the third.}
\end{table}

\subsection{Plain Tables: When NOT to use `ruledtabular'}
There are a number of cases when `ruledtabular' should not
be used: basically for any table using complex content or 
commands.

\subsubsection{Using {\tt\string\multicolumn}}
When you'd like to use the multicolumn command in your
table, you'll find that `ruledtabular' will cause bad
formatting. In that case, Don't Use Ruledtabular, and
instead put in \verb+\hline\hline+ at the top and bottom of
the table. 


\begin{table}[ht]
\caption{A table made without `ruledtabular' needs to have two hlines
added to the top and bottom of the table.}
\vskip3pt
\begin{tabular}{cccccccc}
\hline\hline
 &$r_c$ (\AA)\footnotemark[1]&$r_0$ (\AA)&$\kappa r_0$&
 &$r_c$ (\AA) &$r_0$ (\AA)&$\kappa r_0$\\
\hline
Cu& 0.800 & 14.10 & 2.550 &Sn\footnotemark[1]
& 0.680 & 1.870 & 3.700 \\
Ag& 0.990 & 15.90 & 2.710 &Pb\footnotemark[2]
& 0.450 & 1.930 & 3.760 \\
Au& 1.150 & 15.90 & 2.710 &Ca\footnotemark[3]
& 0.750 & 2.170 & 3.560\\
\hline\hline
\end{tabular}
\footnotetext[1]{This is the first table note.}
\footnotetext[2]{This is the second table note.}
\footnotetext[3]{This is the third table note.}
\end{table}


\subsubsection{Using the {\tt\string\adjustbox\string{\string}
(tabular)\\
\string\end\string{adjustbox\string}} command}
There may be times when the table is too wide, or you want to have
the table be the width of the page, whether or not it appears in
preprint or reprint version of JASA. In this case you can use
\verb+\begin{adjustbox}{<width>}+\\
\verb+(tabular) \end{adjustbox}+.\\
(`adjustbox' will NOT work with `ruledtabular')

You can set a maximum width with \\
\verb+\begin{adjustbox}{max width=\textwidth}(tabular)+\\
\verb+\end{adjustbox}+\\  
in which case the table in the reprint version might be less than
the full text width;

Or you can set the exact width you'd like with\\
\verb+\begin{adjustbox}{width=\textwidth}(tabular)+\\
\verb+\end{adjustbox}+\\
in which case the table will be the full width of the page in
either preprint or reprint.

This way you can make a table that will fit in the correct width
whether you are using the preprint or reprint option.

\subsection{Using dcolumn}
\verb+\usepackage{dcolumn}+ is included in JASA.cls so you
don't need to add it explicitly.
\href{http://anorien.csc.warwick.ac.uk/mirrors/CTAN/macros/latex/required/tools/dcolumn.pdf}
{\url{http://anorien.csc.warwick.ac.uk/mirrors/CTAN/macros/latex/required/tools/dcolumn.pdf}}
will give you detailed information.
A gentler introduction may be found in this informative and well
illustrated
article: \href{https://www.tug.org/pracjourn/2007-1/mori/mori.pdf}
{\url{https://www.tug.org/pracjourn/2007-1/mori/mori.pdf}}, starting on page
20. (You may want to look at more examples in this quite comprehensive
article on making tables in \LaTeX.)

\begin{quote}
``If we do not want to break the fractional and the integral part in two columns,
the dcolumn package provides a new type of column\\
\verb+D{sep -in}{sep -out}{ before.after}+\\
The first argument \verb+{sep-in}+ is the symbol used in the
.tex document to separate
the integral and the fractional part (usually the decimal point . or the decimal
comma ,), the second argument \verb+{sep-out}+
is the symbol that we want in the
output, the third is the number of digits on the left (before) and on the right
(after) this symbol. The numbers are aligned to the decimal point and, in case
that the third argument is negative, the decimal point is aligned to the center of
the column. If the columns have a heading, it must be inserted into
the command \verb+\multicolumn{1}{c}{...}+
\end{quote}


%% Temporarily going into one column mode is very helpful in making
%% up pages with multiple figures and tables. The commands are
%% \onecolumngrid ...\twocolumngrid

\onecolumngrid

\begin{table*}[b]
\caption{Sample using {\tt\string\adjustbox}. Top~5 rated $\widehat{\textrm{ITD}}$ estimation methods
according to the sum and product metric criteria for $\pm0.5$~JND and
$\pm1$~JND tolerance thresholds (normalized scores).}

\begin{adjustbox}{width=\textwidth}
\begin{tabular}{c| l l| l l| l l| l l}
\hline\hline
Rank \#& \multicolumn{2}{c}{sum criteria [$\pm0.5$~JND]} &
\multicolumn{2}{|c}{sum criteria [$\pm1$~JND]} &
\multicolumn{2}{|c}{product criteria [$\pm0.5$~JND]} & \multicolumn{2}{|c}{product criteria [$\pm1$~JND]}\\ 
\hline
1 & Threshold --30dB {lp} & (0.43) & Threshold --30dB {lp} & (0.71) & Threshold --30dB {lp} & (1.00) & Threshold --30dB {lp} & (1.00) \\ 
2 & MaxIACCe {lp} & (0.39) & Threshold --20dB {lp} & (0.66) & MaxIACCe {lp} & (0.39) & Threshold --20dB {lp} & (0.57) \\ 
3 & Threshold --20dB {lp} & (0.38) & CenIACCr {bb} &(0.62) & CenIACCr {lp} &(0.33) & CenIACCr {bb} & (0.37) \\ 
4 & CenIACCr {lp} &(0.37) & MaxIACCe {lp} & (0.61) & Threshold --20dB {lp} & (0.29) & MaxIACCe {lp} & (0.34) \\ 
5 & Cen-$e^2$ {lp} &(0.34) & CenIACCe {lp} & (0.61) & Cen-$e^2$ {lp}
& (0.10) & CenIACCr {lp} & (0.33) \\
\hline\hline
\end{tabular}
\end{adjustbox}

\label{tab:topMethods} 
\end{table*}
\vskip1in
\clearpage
\twocolumngrid

An example using dcolumn:
\begin{verbatim}
{\hsize= 2in
\begin{ruledtabular}
\begin{tabular}{cD {,}{.}{5.4}}
Expression    & \multicolumn {1}{c}{ Value }\\
\hline
$\pi$         &      3,1416             \\
$\pi^{\pi}$   &     36,46               \\
$\pi^{\pi^{\pi}}$ & 80662,7             \\
\end{tabular}
\end{ruledtabular}
}
\end{verbatim}
\vskip12pt
{\hsize= 2in
\begin{ruledtabular}
\begin{tabular}{cD {,}{.}{5.4}}
Expression           & \multicolumn {1}{c}{ Value }\\
\hline
$\pi$                  &      3,1416                 \\
$\pi^{\pi}$           &     36,46                    \\
$\pi^{\pi^{\pi}}$    & 80662,7                      \\
\end{tabular}
\end{ruledtabular}
}

\section{Rotated Table}
On the next page you will see a rotated table.
The commands to rotate your table are\\
\verb+\begin{rotatetable}...\end{rotatetable}+\\
You will not need to use \verb+\begin{table}...\end{table}+
within these commands.

\begin{verbatim}
\begin{rotatetable}
\caption{This is an example of sidewaystable which will
allow a wide table to fit on the page even if it is wider than the 
normal text. We can see where this caption wraps when it is a wide
caption.
}
\begin{tabular}{llllllll}
\hline\hline
 &$r_c$ (\AA)\footnotemark[1]&$r_0$ (\AA)&$\kappa r_0$&
 &$r_c$ (\AA) &$r_0$ (\AA)&$\kappa r_0$\\
\hline
Cu& 0.800 & 14.10 & 2.550 &Sn\footnotemark[1]
& 0.680 & 1.870 & 3.700 \\
Ag& 0.990 & 15.90 & 2.710 &Pb\footnotemark[2]
& 0.450 & 1.930 & 3.760 \\
Au& 1.150 & 15.90 & 2.710 &Ca\footnotemark[3]
& 0.750 & 2.170 & 3.560\\
\hline\hline
Some expressions&make a table&Too wide to fit&in a normal page
width&So we rotate This Table.& We now have the&
height of the page to use& for the table.\\
\end{tabular}
\footnotetext[1]{This is the first table note.}
\footnotetext[2]{This is the second table note.}
\footnotetext[3]{This is the second table note.}
\end{rotatetable}
\end{verbatim}

\begin{rotatetable}
\caption{This is an example of sidewaystable which will
allow a wide table to fit on the page even if it is wider than the 
normal text. We can see where this caption wraps when it is a wide
caption.
}
\begin{tabular}{llllllll}
\hline\hline
 &$r_c$ (\AA)\footnotemark[1]&$r_0$ (\AA)&$\kappa r_0$&
 &$r_c$ (\AA) &$r_0$ (\AA)&$\kappa r_0$\\
\hline
Cu& 0.800 & 14.10 & 2.550 &Sn\footnotemark[1]
& 0.680 & 1.870 & 3.700 \\
Ag& 0.990 & 15.90 & 2.710 &Pb\footnotemark[2]
& 0.450 & 1.930 & 3.760 \\
Au& 1.150 & 15.90 & 2.710 &Ca\footnotemark[3]
& 0.750 & 2.170 & 3.560\\
\hline\hline
Some expressions&make a table&Too wide to fit&in a normal page
width&So we rotate This Table.& We now have the&
height of the page to use& for the table.\\
\end{tabular}
\footnotetext[1]{This is the first table note.}
\footnotetext[2]{This is the second table note.}
\footnotetext[3]{This is the second table note.}
\end{rotatetable}

{
\hsize=\textwidth

 \section{Using Longtable for a table that continues over columns or
pages}
\verb+\usepackage{longtable}+ is included in the JASA.cls, giving you 
access to the commands if you want to make a table that continues
over pages. When using longtable, you must remember to run pdflatex on your article
several times to get correct table formatting.

Here is the syntax for longtable:

\vskip-36pt
{\baselineskip=15pt
\begin{verbatim}
\begin{center} %% Optionally center table with \begin{center} ...\end{center}
\begin{longtable}{<table preamble>}
\caption{} %% optional caption
   %% Everything between here and \endfirsthead will be used for
   %% column headers for the first page of the table. Use two \hline's here and 
   %% at the bottom of the table.
\hline\hline \multicolumn{1}{|c|}{\textbf{First column}} &
\multicolumn{1}{c|}{\textbf{Second column}}  &
\multicolumn{1}{c|}{\textbf{Third column}} \\ \hline 
\endfirsthead
   %% Everything between here and \endhead will be used for
   %% column headers for the all the following pages of the table.
\multicolumn{3}{c}%
{{\bfseries \tablename\ \thetable{} -- continued from previous page}} \\
\hline \multicolumn{1}{|c|}{\textbf{First column}} & 
\multicolumn{1}{c|}{\textbf{Second column}} & 
\multicolumn{1}{c|}{\textbf{Third column}} \\ \hline 
\endhead
   %% Everything between here and \endfoot will be used for footers on every
   %% page of the table except for the final one:
\hline \multicolumn{3}{|r|}{{Continued on next page}} \\ \hline
\endfoot
   %% Everything between here and \endlastfoot will be used for 
   %% bottom of the table on its final page:
\hline \hline
\endlastfoot
   %% Enter contents of the table here:
Table Text
   %% Then end table with
\end{longtable}
\end{center}
\end{verbatim}
}}


\clearpage
An example Long Table:
\begin{center}
\tabcolsep=12pt
\begin{longtable}{|l|l|l|}
\caption{A sample long table.} \label{tab:long} \\
\hline\hline \multicolumn{1}{|c|}{\textbf{First}} &
\multicolumn{1}{c|}{\textbf{Second}}  &
\multicolumn{1}{c|}{\textbf{Third}} \\ \hline 
\endfirsthead
    \multicolumn{3}{c}
    {\tablename\ \thetable\ -- \textit{Continued}} \\
\hline \multicolumn{1}{|c|}{\textbf{\itshape First}} &
\multicolumn{1}{c|}{\textbf{\itshape Second }}  &
\multicolumn{1}{c|}{\textbf{\itshape Third }} \\ \hline 
\endhead
\hline
\multicolumn{3}{l}{\it(Continued)}
\endfoot
\hline \hline
\endlastfoot
And&So&On\\
And&So&On\\
And&So&On\\
And&So&On\\
And&So&On\\
And&So&On\\
And&So&On\\
And&So&On\\
And&So&On\\
And&So&On\\
And&So&On\\
And&So&On\\
And&So&On\\
And&So&On\\
And&So&On\\
And&So&On\\
And&So&On\\
And&So&On\\
And&So&On\\
And&So&On\\
And&So&On\\
And&So&On\\
And&So&On\\
And&So&On\\
And&So&On\\
And&So&On\\
And&So&On\\
And&So&On\\
And&So&On\\
And&So&On\\
And&So&On\\
And&So&On\\
And&So&On\\
And&So&On\\
And&So&On\\
And&So&On\\
And&So&On\\
And&So&On\\
And&So&On\\
And&So&On\\
And&So&On\\
And&So&On\\
And&So&On\\
And&So&On\\
And&So&On\\
\end{longtable}
\end{center}

\subsection{A long table with rotated column heads}
This example shows how to rotate column heads in a long table
environment. You may want to try these commands on another
table where the column headers are too wide to fit gracefully
in the column or page.

To rotate the column header use the command \verb+\rothead{<column head>}+.

In this example, the code for the rotated columns heads looks like
this:

\begin{verbatim}
\rothead{Bin Interval\\ Lower Bound} & 
\rothead{Bin Interval\\ Upper Bound} & 
\rothead{Histogram\\ Count} & 
\rothead{Expected\\ Count} & 
\rothead{Cumulative\\ Distribution} & 
\rothead{$\chi^2$-Value} \\
\end{verbatim}

%%%%%%%%%%%
%% For the table following to make column heads rotate 90 degrees.
\renewcommand{\theadfont}{\normalsize\boldmath\bfseries}
\settowidth{\rotheadsize}{\bfseries Upper Bound}
\setlength{\extrarowheight}{2pt}

\begin{longtable}{|c|c|c|c|c|c|}
 \caption{Longtable sample, table that continues over pages, and
in this case, rotates column heads.}
\label{Table1}\\
    \hline\hline
    \rothead{Bin Interval\\ Lower Bound} & \rothead{Bin Interval\\ Upper Bound} & \rothead{Histogram\\ Count} & \rothead{Expected\\ Count} & \rothead{Cumulative\\ Distribution} & \rothead{$\chi^2$-Value} \\
    \hline
    \endfirsthead
    \multicolumn{6}{c}
    {\tablename\ \thetable\ -- \textit{Continued from previous page or
    column}} \\
    \hline
    \rothead{Bin Interval\\ Lower Bound} & \rothead{Bin Interval\\ Upper Bound} & \rothead{Histogram\\ Count} & \rothead{Expected\\ Count} & \rothead{Cumulative\\ Distribution} & \rothead{$\chi^2$-Value} \\
    \hline
    \endhead
    \hline \multicolumn{6}{r}{\textit{Continued on following page or
    column}} \\
    \endfoot
    \hline\hline
    \endlastfoot
        1000 & 1003,99 & 102 & 100 & 102 & 0,04 \\
        1004 & 1007,99 & 105 & 100 & 207 & 0,25 \\
        1008 & 1011,99 & 104 & 100 & 311 & 0,16 \\
        1012 & 1015,99 & 104 & 100 & 415 & 0,16 \\
        1000 & 1003,99 & 102 & 100 & 102 & 0,04 \\
        1004 & 1007,99 & 105 & 100 & 207 & 0,25 \\
        1008 & 1011,99 & 104 & 100 & 311 & 0,16 \\
        1012 & 1015,99 & 104 & 100 & 415 & 0,16 \\
        1000 & 1003,99 & 102 & 100 & 102 & 0,04 \\
        1004 & 1007,99 & 105 & 100 & 207 & 0,25 \\
        1008 & 1011,99 & 104 & 100 & 311 & 0,16 \\
        1012 & 1015,99 & 104 & 100 & 415 & 0,16 \\
        1000 & 1003,99 & 102 & 100 & 102 & 0,04 \\
       1004 & 1007,99 & 105 & 100 & 207 & 0,25 \\
\end{longtable}

\begin{figure*}[p]
\parskip=6pt
\baselineskip=12pt

\subsection{Sample Figures, new commands available in this style}

{\it Note that the publisher determines the final layout, so 
your choice of figure alignment may not be reflected in the published
article.}

\noindent
\verb+\figline{}+ will center one or more figures on one line. 


\noindent
\verb+\fig{<name of file>}{<width>}{<letter to put underneath>}+

If you {\bf do not want to use a lettered label} you should use facing
curly brackets for the third argument of \verb+\fig+:

\verb+\figline{\fig{<name of file>}{<width>}{}}+

\noindent
\verb+\leftfig{<name of file>}{<width>}{<letter to put underneath>}+

\noindent
\verb+\rightfig{<name of file>}{<width>}{<letter to put underneath>}+

\noindent
\verb+\boxedfig{<name of file>}{<width>}{<letter to put underneath>}+

\noindent
\verb+\rotatefig{<degrees of rotation>}{<name of file>}{<width>}+\\
\verb+{<letter to put underneath>}+

The following illustrations show these commands in use.


\begin{verbatim}
\figline{\fig{figsamp}{4cm}{(a)}
\fig{figsamp}{4cm}{(b)}}
\figline{\fig{figsamp}{4cm}{(c)}
\fig{figsamp}{4cm}{(d)}}
\figline{\fig{figsamp}{4cm}{(e)}}
\end{verbatim}

\figline{\fig{figsamp}{4cm}{(a)}
\fig{figsamp}{4cm}{(b)}}
\figline{\fig{figsamp}{4cm}{(c)}
\fig{figsamp}{4cm}{(d)}}
\figline{\fig{figsamp}{4cm}{(e)}}

\caption{ \label{fig:pressure_field} Multiple images on one figure example (a) image 1, (b-f) ($\rho$=1000 kg/m$^3$) and speed of sound ($c$=1500 m/s).}

\end{figure*}

\begin{figure*}
\baselineskip=12pt
\begin{verbatim}
\figline{\boxedfig{figsamp}{2in}{(a)}}
\figline{\leftfig{figsamp}{2in}{(b)}\rightfig{figsamp}{2in}{(c)}}
\figline{\rotatefig{90}{figsamp}{2in}{(d)}\rotatefig{180}{figsamp}{2in}{(e)}}
\end{verbatim}

\figline{\boxedfig{figsamp}{2in}{(a)}}

\figline{\leftfig{figsamp}{2in}{(b)}\rightfig{figsamp}{2in}{(c)}}

\figline{\rotatefig{90}{figsamp}{2in}{(d)}
\rotatefig{180}{figsamp}{2in}{(e)}
}

\caption{More figure examples: (a) boxedfig, 
(b)leftfig; (c)right fig; (d) rotatefig 90 degrees;
(e) rotatefig 180 degrees. }

\end{figure*}

\begin{figure*}[h]
\section{Labels in Figline}
We can label and
reference separate parts of the figure when using figline.
The reference will give the illustration letter as well as the figure
number.

To label figures used in \verb+\figline{}+ type in your
label immediately after the \verb+\fig{}{}{}+ command, inside the
argument to figline. 

\begin{verbatim}
\figline{\fig{<name of file>}{<width>}{<letter to put underneath>}\label{<labelname>}}
\end{verbatim}


The same placement should be used for all the kinds of fig
environments used in \verb+\figline{}+:
\begin{verbatim}
\fig{}{}{}\label{}, \leftfig{}{}{}\label{}, \rightfig{}{}{}\label{}, 
\boxedfig{}{}{}\label{}, \rotatefig{}{}{}\label{}  
\end{verbatim}

\baselineskip=12pt
\begin{verbatim}
\figline{\boxedfig{figsamp}{1in}{(a)}\label{boxedfigLetter}}
\figline{\leftfig{figsamp}{1in}{(b)}
\rightfig{figsamp}{1in}{(c)}\label{rightfigLetter}}
\figline{\rotatefig{90}{figsamp}{1in}{(d)}
\rotatefig{180}{figsamp}{1in}{(e)}\label{rotatefigLetter}}
\end{verbatim}

\figline{\boxedfig{figsamp}{1.5in}{(a)}\label{boxedfigLetter}}
\figline{\leftfig{figsamp}{1.5in}{(b)}\rightfig{figsamp}{1.5in}{(c)}\label{rightfigLetter}}
\figline{\rotatefig{90}{figsamp}{1.5in}{(d)}\rotatefig{180}{figsamp}{1.5in}{(e)}\label{rotatefigLetter}}

\caption{More figure examples, showing how to enter
{\tt\string\label\string{\string}} command.}


\verb+References:\ref{boxedfigLetter},\ref{rightfigLetter},\ref{rotatefigLetter}+ 

which produces\\

References: \ref{boxedfigLetter},\ref{rightfigLetter},\ref{rotatefigLetter}


\end{figure*}

\begin{figure*}
\baselineskip=12pt
\begin{verbatim}
\sidebysidefigures{figsamp}{Describing the first
illustration.}/{figsamp}{Describing the second illustration.}
\end{verbatim}

\sidebysidefigures{figsamp}{Describing the first
illustration.}/{figsamp}{Describing the second illustration.}

\baselineskip=12pt
\begin{verbatim}
\figline{\fig{figsamp}{.2\textwidth}{(A)}
\fig{figsamp}{.2\textwidth}{(B)}
\fig{figsamp}{.2\textwidth}{(C)}}
\caption{Caption for three illustrations. 
The caption may produce many lines, but only one paragraph.
}
\end{verbatim}
\figline{\fig{figsamp}{.2\textwidth}{(A)}
\fig{figsamp}{.2\textwidth}{(B)}
\fig{figsamp}{.2\textwidth}{(C)}}
\caption{Caption for three illustrations. 
The caption may produce many lines, but only one paragraph.  The caption may produce many lines, but only one paragraph.
}
\end{figure*}

\begin{figure*}
\baselineskip=12pt
\begin{verbatim}
\figline{\hfill\fig{figsamp}{.2\textwidth}{(a)}
\fig{figsamp}{.2\textwidth}{(b)}\hfill}
\figline{\hfill\fig{figsamp}{.2\textwidth}{(c)}
\fig{figsamp}{.2\textwidth}{(d)}\label{test}\hfill}
\caption{Here is a caption.
You cannot have more than one paragraph of text in a caption.
You cannot have more than one paragraph of text in a caption.
You cannot have more than one paragraph of text in a caption.
You cannot have more than one paragraph of text in a caption.
}
\end{verbatim}
\figline{\hfill\fig{figsamp}{.2\textwidth}{(a)}
\fig{figsamp}{.2\textwidth}{(b)}\hfill}
\figline{\hfill\fig{figsamp}{.2\textwidth}{(c)}
\fig{figsamp}{.2\textwidth}{(d)}\label{test}\hfill}
\caption{Here is a caption.
You cannot have more than one paragraph of text in a caption.
You cannot have more than one paragraph of text in a caption.
You cannot have more than one paragraph of text in a caption.
You cannot have more than one paragraph of text in a caption.
}
\vskip1in
\end{figure*}
\clearpage
\begin{figure}[h]
\baselineskip=12pt
\begin{verbatim}
\figcolumn{
\fig{figsamp}{.2\textwidth}{(A)}
\fig{figsamp}{.2\textwidth}{(B)}
\fig{figsamp}{.2\textwidth}{(C)}
}
\end{verbatim}

\figcolumn{
\fig{figsamp}{.2\textwidth}{(A)}
\fig{figsamp}{.2\textwidth}{(B)}
\fig{figsamp}{.2\textwidth}{(C)}
}

\caption{Here are some stacking figures in a single column. 
The {\tt\string\figcolumn\string{\string}} command works equally well in single
or double column figures.}
\end{figure}

\begin{nofloatfigure}
%% \reprintcolumnwidth is the same in preprint and reprint for
%% ease of use for authors:
\includegraphics[width=\reprintcolumnwidth]{figsamp}
\caption{\label{fig:FIG1}
Here is a sample of a long caption that will continue
over columns or pages because it is used in 
{\tt\string\begin\string{nofloatfigure\string}...\string\end\string{nofloatfigure\string}}
We want to test a very long caption, so here it is.
We want to test a very long caption, so here it is.
We want to test a very long caption, so here it is.
We want to test a very long caption, so here it is.
We want to test a very long caption, so here it is.
We want to test a very long caption, so here it is.
We want to test a very long caption, so here it is.
We want to test a very long caption, so here it is.
We want to test a very long caption, so here it is.
We want to test a very long caption, so here it is.
}
\end{nofloatfigure}

\section{Algorithm Example}
JASA.cls includes\\ \verb+\usepackage{algorithm2e}+,\\
\verb+\usepackage{algorithmic}+, and\\ \verb+\usepackage{algcompatible}+.

Below is an example of\\
\verb+\begin{algorithmic}...\end{algorithmic}+\\
used within\\
\verb+\begin{algorithm}...\end{algorithm}+.\\
Note that the commands
that are printed in bold are all entered with all caps.

\begin{algorithm}[ht]
\caption{Sample code is shown using the algorithmic commands without
numbering.}
\begin{algorithmic}
\IF {$i\geq maxval$}
    \STATE $i\gets 0$
\ELSE
    \IF {$i+k\leq maxval$}
        \STATE $i\gets i+k$
    \ENDIF
\ENDIF
\end{algorithmic}
\end{algorithm}

\begin{algorithm}[ht]
\caption{Sample code is shown using the algorithmic commands {\it
with}
numbering, by following algorithmic with {\tt [1]}, ie,
({\tt\string\begin\string{algorithmic\string}[1]}).}
\begin{algorithmic}[1]
\IF {$i\geq maxval$}
    \STATE $i\gets 0$
\ELSE
    \IF {$i+k\leq maxval$}
        \STATE $i\gets i+k$
    \ENDIF
\ENDIF
\end{algorithmic}
\end{algorithm}

Documentation for the algorithm2e  commands:\\
\url{http://tug.ctan.org/macros/latex/contrib/algorithm2e/doc/algorithm2e.pdf}

Documentation for the algorithmicx  commands:\\
\url{http://tug.ctan.org/macros/latex/contrib/algorithmicx/algorithmicx.pdf}

A description of options for the algorithm bundle found here:\\
\url{http://ctan.math.utah.edu/ctan/tex-archive/macros/latex/contrib/algorithms/algorithms.pdf}


\subsection{Example of multimedia entry}
Please note that this is for multimedia intended to appear inline
within the published article. 

Here is what a multimedia entry will look like:
\multimedia{http://dx.doi.org/10.1121/1.4947423.1}{Corresponding pulse-compressed echo envelopes
and video recordings from a fluttering luna moth.
Echoes from the wings and body of the moth generally dominate the
acoustic returns, which vary greatly over consecutive ensonifications
across the wingbeat cycle. File of type ``mp4'' (15.3
MB)}\label{mmtest1}

Here we try cross referencing the multimedia entry: The multimedia
above is Mm.~\ref{mmtest1}.

\subsection{Supplementary Material}
ASA
prefers that authors to submit related/relevant article files as
supplementary material with their submission.

\subsection{Supplementary material for publication}
Any archival supplemental materials to be published with the
manuscript (eg., supplementary figures) should be cited in-text and a footnote provided.

An example of reference to supplementary material:

The sound files and videos for this and other figures
are included as supplementary materials\footnote{See
Supplementary materials at [URL will be inserted by AIP]
for [give a brief description of the material].}.

The contents of the footnote above will appear at the beginning of the
bibliography made with BibTeX when the default `author-year' documentclass option is used;
BibTeX output will have the footnote interleaved with other
references if the NumberedRefs documentclass option is used.

\subsection{File naming conventions}
Here are the conventions for naming files:

\begin{itemize}
\item
Supplementary Figure or
	Supplementary Figure or Text files should be named: SuppPub\#.xxx, where ``\#'' is
	a number and ``xxx" is the file format extension
	(SuppPub1.docx, SuppPub2.jpg, etc)

\item
	Supplementary Multimedia files: SuppPubmm\#.xxx, where ``\#'' is a
	number and ``xxx'' is the file format extension (SuppPubmm1.mp3,
	SuppPubmm2.gif, etc)

\item
Multimedia files must be named accordingly: MM\#.xxx, where ``\#'' is the
number and ``xxx'' is the file format extension (MM1.wav, MM2.avi, etc).

\item
The only figure formats allowed are the following: 
.pdf, .ps, .eps, or .jpg. Figure files must be named in this fashion:
Figure\#.xxx, where ``\#'' is the figure number and ``xxx'' is the file format
(Figure1.eps, Figure2.jpg, Figure3a.ps, Figure3b.ps, etc). 

\end{itemize}



\section{\label{sec:5}Conclusion}

And in conclusion\ldots

\begin{acknowledgments}
This research was supported by  ...
\end{acknowledgments}


 % -------------------------------------------------------------------------------------------------------------------
 %   Appendix  (optional)

\section{Appendices}

To start the appendix, use the \verb+\appendix+ command.
This signals that all following section commands will produce appendixes
instead of regular sections. 

Therefore, the \verb+\appendix+ command
should be used only once---to set up the section commands to act as
appendices. Subsection and subsubsections are not changed from normal subsection
and subsubsection commands.

\docsection{Making Multiple Appendices}
Every \verb+\section+ command after \verb+\appendix+ will produce a new appendix with 
a new appendix letter:

\appendix
\section{One of multiple appendices}
\subsection{Here is subsection}
\subsubsection{Here is subsubsection}


\section{another appendix}
Notice that the equation will use the letter for the current appendix:

\begin{equation}
t=-\frac{1}{\bar{\chi}_\beta}\frac{\sin^2\phi}{\cos\bar\theta}+t_h
\end{equation}

\docsection{Appendix without a title}
To make an appendix without a title, after \verb+\appendix+ use \verb+\section{}+. For example,

\begin{verbatim}
\appendix
\section{}
\end{verbatim}
will produce:

\appendix
\section{}

\newpage
\docsection{Subsection or subsubsection in an appendix}

\subsection{A subsection in an appendix}
\label{app:subsec}
Note the equation numbers in this appendix, produced with the
subequations environment:
\begin{subequations}
\begin{eqnarray}
E&=&mc, \label{appa}
\\
E&=&mc^2, \label{appb}
\\
E&\agt& mc^3. \label{appc}
\end{eqnarray}
\end{subequations}

\subsubsection{A subsubsection in an appendix}
References for subequations:
they turn out to be Eqs.~(\ref{appa}), (\ref{appb}), and (\ref{appc}).

\subsection{Labels should go AFTER appendix title}
This will work:
\begin{verbatim}
\section{Testing}
\label{app:testing}
\end{verbatim}

This will not:
\begin{verbatim}
\section{Testing\label{app:testing}}
\end{verbatim}


\docsection{Only one appendix? use {\tt \string\appendix*}}
If there is only one appendix, then the letter ``A'' should not
appear. This is suppressed by using the star version of the appendix
command (\verb+\appendix*+ in the place of \verb+\appendix+).

Since this is a single appendix, the {\bf first} \verb+\section{}+ command 
after \verb+\appendix*+ will make an appendix heading. 
However, after the first section, \verb+\section+ 
{\bf will not} produce the word `Appendix', but will
be a plain section head.

Here are examples of single appendices that are not lettered.

\docsection{Single appendix with title}

\appendix*
%% only needed because we have several \appendix* commands
%% in the same document:


\section{Single Appendix}
\subsection{Here is subsection}
\subsubsection{Here is subsubsection}

All of the equations in the single appendix will use the letter `A':
\begin{equation}
t=-\frac{1}{\bar{\chi}_\beta}\frac{\sin^2\phi}{\cos\bar\theta}+t_h
\end{equation}

\newpage


\docsection{Single appendix without title}

\appendix*

\section{}
\subsection{Here is the topic of this appendix}



\docsection{Sample appendix figure and table}
Figure and table numbering are continuous through the article,
and handled the same as they are in the rest of the article.

\begin{figure}[ht]
\includegraphics[width=1in]{figsamp}
\caption{Figure in an appendix.}
\end{figure}


\begin{table}[ht]
\caption{Here is the caption for a table in an appendix.}
\centering
\begin{ruledtabular}
\begin{tabular}{cccc}
one&two&three&four\\
\hline
C&D&E&F\\
\end{tabular}
\end{ruledtabular}
\vskip12pt
\end{table}


\docsection{Footnotes}
The contents of the footnotes will appear at the beginning of the
bibliography when BibTeX produces the .bbl file using the default
AuthorYear style; interleaved with other references if 
Numbered\-Refs option:\\
\verb+\documentclass[preprint,NumberedRefs]{JASA}+\\
and BibTeX has been used.

This example show where this cite \citep{booksamp1} will appear in the
bibliography,
depending on whether we use default author-year style
or call for the NumberedRefs documentclass option.

This example shows what happens when there are two references to the
same author and year,  \citet{Shera2} and \citet{Shera}.

Here are some sample footnotes:\footnote{Here is the second footnote.
It will appear before the beginning of the bibliography in Author-Year
style (default) or it will be 
 interleaved with other references when using the Numbered\-Refs
 option.}$^,$% (comma between two sequential footnotes)
\footnote{Here is a third footnote.}


\docsection{Making the Bibliography Using BibTeX}
Authors are highly  recommended to use BibTeX to produce their
bibliographies. The results will be predictable and even if
it might take some time to get comfortable with  using BibTeX,
in the long run it will save you endless aggravation.

A resource for making your bibliography entries
correctly is included in this package: 
ReferenceStyles.pdf. You will also find
the files
bibsamp1.tex/.pdf and bibsamp2.tex/.pdf
for examples of output; and sampbib.bib for an example of
how to make your .bib database entries.

There are two possible bibliography styles: the default, author-year,
and the optional style, Numbered\-Refs, which you would call using\\
{\verb+\documentclass[reprint,NumberedRefs]{JASA}+ }

\verb+\citep{}+ should normally be used rather than \verb+\cite{}+.



You can also use \verb+\citet{}+ if it is more grammatically correct
to have only the year in parens (note: this is used with author-year style references).

\begin{verbatim}
\citep{bibitemName} = (bibitemName, year)
or
\citet{bibitemName} = bibitemName (year)
\end{verbatim}

Note that the citations are hyperlinked to their entries in the
bibliography:

Normal journal cite: \citep{joursamp1},
 Book reference \citet{booksamp1}, 
Computer language documentation:
\citep{sampcode2}.

Every \verb+\citep+  or \verb+\citet{}+ will produce a citation and an entry in the
bibliography. Every citation must have a matching entry in the
bibliography
database file (\verb+\filename.bib+).

Make your bibliography by doing: pdflatex filename,  bibtex filename,
pdflatex filename, pdflatex filename.


When uploading your files to Editorial Manager,  include both the
.bib and the appropriate .bst file (for\break 
\newpage
\noindent
author/year reference style: jasaauthoryear2.bst; for numerical style: jasanum2.bst). Both the .bib and .bst should be uploaded as the “Manuscript (TeX or Word only)” item type.


Compare the results you get with\\
\verb+\documentclass[reprint]{JASA}+ \\
vs.\\
\verb+\documentclass[reprint,NumberedRefs]{JASA}+ 


\bibliography{sampbib}

\end{document}
